\chapter{Wdrożenie i pomiary kosztów w sieciach warstwy drugiej}
\label{ch:l2}

Rozdział~\ref{ch:gaz} wykazał, że po stronie samego kontraktu koszt publikacji
korzenia Merkle jest niemal stały i sprowadza się do pojedynczej, taniej operacji
zapisu. Rzeczywisty koszt ponoszony przez uczelnię zależy jednak nie tylko od
zużycia gazu, lecz przede wszystkim od sieci, w~której kontrakt zostaje wdrożony.
Warstwa pierwsza Ethereum zapewnia najsilniejsze gwarancje bezpieczeństwa, ale jej
koszt bywa o~rzędy wielkości wyższy niż w~sieciach warstwy drugiej (ang.
\emph{layer~2}, L2), które przenoszą wykonanie poza łańcuch główny, publikując na
nim jedynie skompresowane dane. Niniejszy rozdział mierzy ten koszt empirycznie:
wdraża identyczny kontrakt \code{DiplomaRegistry} w~czterech sieciach i porównuje
koszt kluczowych operacji zarówno w~jednostkach gazu, jak i w~dolarach
amerykańskich.

\section{Metodyka pomiaru}
\label{sec:l2-metodyka}

\subsection{Zakres i badane sieci}
Porównanie obejmuje koszt uruchomienia kontraktu \code{DiplomaRegistry} na czterech
sieciach Ethereum: warstwie pierwszej (Sepolia, jako punkt odniesienia) oraz trzech
rozwiązaniach warstwy drugiej reprezentujących odmienne rodziny technologiczne ---
Arbitrum Nitro, Base jako przedstawiciel stosu OP oraz zkSync~Era jako rollup
oparty na dowodach zwięzłości (ang. \emph{validity proofs}). Dobór ten pozwala
porównać nie pojedyncze wdrożenia, lecz reprezentatywne klasy rollupów.

\subsection{Pomiary}
Każdą sieć mierzono na jej publicznej sieci testowej. \emph{Zużycie gazu} oraz
\emph{opłata za dane L1} są wielkościami rzeczywistymi, odczytanymi z~paragonów
transakcji, natomiast koszty w~USD modelowane są na podstawie historycznych
wartości \code{baseFeePerGas} warstwy pierwszej z~ostatnich 90~dni (percentyle
p10/p50/p90). Wszystkie wyniki w~USD przeliczono kursem
$\text{ETH}/\text{USD}=3500$ (snapshot z~dnia 13~maja~2026~r.); wpływ tego założenia
badany jest w~sekcji~\ref{sec:l2-wrazliwosc}. Pomiary wykonano 19~maja~2026~r.

\subsection{Parametry próby}
Dla operacji \code{issueBatch} użyto rozmiarów paczek
$\{1, 10, 100, 1\,000, 10\,000\}$. Dla \code{revokeFromBatch} zebrano 30~próbek
przeciwko paczce bazowej o~rozmiarze 40, z~losowym wyborem unieważnianych liści.
Dla \code{statusWithProof} zebrano 100~próbek czasu odpowiedzi RPC. Każdy pomiar
opóźnienia powtórzono $N{=}30$ razy, a~raportowane wartości są średnią z~trzech
niezależnych przebiegów pomiarowych.

\section{Koszt operacji \code{issueBatch}}
\label{sec:l2-issue}

Operacja \code{issueBatch} publikuje korzeń Merkle całej paczki dyplomów.
Tabela~\ref{tab:issue-gas} pokazuje jej zużycie gazu, a~tabela~\ref{tab:issue-usd}
--- odpowiadający mu koszt w~USD. Ponieważ na łańcuchu zapisywany jest zawsze
pojedynczy 32-bajtowy korzeń, zużycie gazu jest praktycznie niezależne od liczby
dyplomów w~paczce; różnicę kosztu między sieciami tworzy sposób rozliczania danych,
nie rozmiar danych.

\begin{table}[H]
  \centering
  \caption{Zużycie gazu operacji \code{issueBatch} według sieci i~rozmiaru paczki
  (wartości deterministyczne).}
  \label{tab:issue-gas}
  \pgfplotstabletypeset[
    col sep=comma, string type,
    every head row/.style={before row=\toprule, after row=\midrule},
    every last row/.style={after row=\bottomrule},
    columns={chain,gas_1,gas_10,gas_100,gas_1000,gas_10000},
    columns/chain/.style={column name=Sieć, column type=l},
    columns/gas_1/.style={column name={$n{=}1$}, column type=r},
    columns/gas_10/.style={column name={$n{=}10$}, column type=r},
    columns/gas_100/.style={column name={$n{=}100$}, column type=r},
    columns/gas_1000/.style={column name={$n{=}1000$}, column type=r},
    columns/gas_10000/.style={column name={$n{=}10000$}, column type=r},
  ]{\DataDir/table-issue-cost.csv}
\end{table}

\begin{table}[H]
  \centering
  \caption{Koszt operacji \code{issueBatch} w~USD (scenariusz p50,
  $\text{ETH}/\text{USD}{=}3500$). Pełną zmienność scenariuszy p10/p50/p90 dla
  $n{=}1000$ przedstawia rys.~\ref{fig:basefee}.}
  \label{tab:issue-usd}
  \pgfplotstabletypeset[
    col sep=comma, string type,
    every head row/.style={before row=\toprule, after row=\midrule},
    every last row/.style={after row=\bottomrule},
    columns={chain,usd_p50_1,usd_p50_10,usd_p50_100,usd_p50_1000,usd_p50_10000},
    columns/chain/.style={column name=Sieć, column type=l},
    columns/usd_p50_1/.style={column name={$n{=}1$}, column type=r},
    columns/usd_p50_10/.style={column name={$n{=}10$}, column type=r},
    columns/usd_p50_100/.style={column name={$n{=}100$}, column type=r},
    columns/usd_p50_1000/.style={column name={$n{=}1000$}, column type=r},
    columns/usd_p50_10000/.style={column name={$n{=}10000$}, column type=r},
  ]{\DataDir/table-issue-cost.csv}
\end{table}

Praktyczną miarą jest koszt w~przeliczeniu na jeden dyplom, malejący wraz ze
wzrostem paczki dzięki amortyzacji stałego kosztu publikacji korzenia
(tabela~\ref{tab:per-diploma}, rys.~\ref{fig:cost-per-diploma}).

\begin{table}[H]
  \centering
  \caption{Koszt na dyplom (amortyzacja paczki) w~USD, scenariusz p50.}
  \label{tab:per-diploma}
  \pgfplotstabletypeset[
    col sep=comma, string type,
    every head row/.style={before row=\toprule, after row=\midrule},
    every last row/.style={after row=\bottomrule},
    columns/chain/.style={column name=Sieć, column type=l},
    columns/usd_p50_per_diploma_1/.style={column name={$n{=}1$}, column type=r},
    columns/usd_p50_per_diploma_10/.style={column name={$n{=}10$}, column type=r},
    columns/usd_p50_per_diploma_100/.style={column name={$n{=}100$}, column type=r},
    columns/usd_p50_per_diploma_1000/.style={column name={$n{=}1000$}, column type=r},
    columns/usd_p50_per_diploma_10000/.style={column name={$n{=}10000$}, column type=r},
  ]{\DataDir/table-issue-per-diploma.csv}
\end{table}

\begin{figure}[H]
  \centering
  \begin{tikzpicture}
    \begin{loglogaxis}[
      xlabel={Rozmiar paczki (liczba dyplomów)},
      ylabel={Koszt na dyplom (USD)},
      legend pos=north east, grid=major,
      width=0.85\textwidth, height=8cm,
    ]
      \addplot table[x=batchSize, y=sepolia] {\DataDir/fig1-cost-per-diploma.dat};
      \addplot table[x=batchSize, y=arbitrumSepolia] {\DataDir/fig1-cost-per-diploma.dat};
      \addplot table[x=batchSize, y=baseSepolia] {\DataDir/fig1-cost-per-diploma.dat};
      \addplot table[x=batchSize, y=zksyncSepolia] {\DataDir/fig1-cost-per-diploma.dat};
      \legend{Sepolia L1, Arbitrum Sepolia, Base Sepolia, zkSync Sepolia}
    \end{loglogaxis}
  \end{tikzpicture}
  \caption{Koszt na dyplom w~funkcji rozmiaru paczki (skala log-log, scenariusz
  p50).}
  \label{fig:cost-per-diploma}
\end{figure}

\section{Koszt operacji \code{revokeFromBatch}}
\label{sec:l2-revoke}

Unieważnienie pojedynczego dyplomu (\code{revokeFromBatch}) jest operacją rzadką i
dotyczy jednej transakcji. Tabela~\ref{tab:revoke} podaje jej medianowy koszt;
utrzymuje on ten sam porządek sieci co koszt publikacji, z~Base jako
zdecydowanie najtańszą opcją.

\begin{table}[H]
  \centering
  \caption{Koszt operacji \code{revokeFromBatch} (pojedyncza transakcja, mediana).}
  \label{tab:revoke}
  \pgfplotstabletypeset[
    col sep=comma, string type,
    every head row/.style={before row=\toprule, after row=\midrule},
    every last row/.style={after row=\bottomrule},
    columns={chain,median_gas,median_l1data_wei,usd_p50},
    columns/chain/.style={column name=Sieć, column type=l},
    columns/median_gas/.style={column name={Gaz (mediana)}, column type=r},
    columns/median_l1data_wei/.style={column name={Dane L1 [wei]}, column type=r},
    columns/usd_p50/.style={column name={Koszt [USD]}, column type=r},
  ]{\DataDir/table-revoke-cost.csv}
\end{table}

\section{Latencja}
\label{sec:l2-latencja}

Poza kosztem istotny jest czas odpowiedzi. Tabela~\ref{tab:read-latency} podaje
latencję odczytu \code{statusWithProof} (weryfikacja po stronie sprawdzającego),
a~tabela~\ref{tab:inclusion-latency} --- czas od wysłania transakcji zapisu do
odebrania paragonu przez klienta.

\begin{table}[H]
  \centering
  \caption{Latencja odczytu \code{statusWithProof} (ms).}
  \label{tab:read-latency}
  \pgfplotstabletypeset[
    col sep=comma,
    every head row/.style={before row=\toprule, after row=\midrule},
    every last row/.style={after row=\bottomrule},
    columns={chain,p50_ms,p95_ms,max_observed_ms},
    columns/chain/.style={string type, column name=Sieć, column type=l},
    columns/p50_ms/.style={column name={$p_{50}$}, column type=r, fixed, fixed zerofill, precision=1},
    columns/p95_ms/.style={column name={$p_{95}$}, column type=r, fixed, fixed zerofill, precision=1},
    columns/max_observed_ms/.style={column name={maks.}, column type=r, fixed, fixed zerofill, precision=1},
  ]{\DataDir/table-read-latency.csv}
\end{table}

\begin{table}[H]
  \centering
  \caption{Czas odbioru paragonu transakcji zapisu po stronie klienta (ms).
  Kolumny \code{issue} pochodzą z~$N{=}30$ powtórzeń \code{issueBatch(1000)}
  w~każdym przebiegu i~są agregowane po wszystkich trzech przebiegach.}
  \label{tab:inclusion-latency}
  \pgfplotstabletypeset[
    col sep=comma, string type,
    every head row/.style={before row=\toprule, after row=\midrule},
    every last row/.style={after row=\bottomrule},
    columns/chain/.style={column name=Sieć, column type=l},
    columns/revoke_p50_ms/.style={column name={revoke $p_{50}$}, column type=r},
    columns/revoke_p95_ms/.style={column name={revoke $p_{95}$}, column type=r},
    columns/issue_p50_ms/.style={column name={issue $p_{50}$}, column type=r},
    columns/issue_p95_ms/.style={column name={issue $p_{95}$}, column type=r},
    columns/issue_sigma_ms/.style={column name={issue $\sigma$}, column type=r},
    columns/issue_n/.style={column name={$N$}, column type=r},
  ]{\DataDir/table-inclusion-latency.csv}
\end{table}

Mierzoną wielkością jest czas między wywołaniem \code{sendTransaction} a~otrzymaniem
paragonu przez klienta RPC. Przy interwale odpytywania około 4~s (domyślnym
w~bibliotece viem) wynik jest górnym ograniczeniem faktycznego czasu inkluzji
w~bloku, nie zaś jego dokładnym pomiarem; rozkład wartości na sieci Sepolia jest
wyraźnie skupiony wokół wielokrotności czasu bloku ($\approx 12$~s), co potwierdza
tę interpretację.

\section{Analiza wrażliwości}
\label{sec:l2-wrazliwosc}

Ponieważ koszty w~USD zależą od trzech zmiennych rynkowych --- opłaty bazowej L1,
kursu ETH/USD oraz ceny gazu sekwensera L2 --- poniżej badana jest odporność
wyników na ich wahania.

\subsection{Scenariusze opłaty bazowej L1}
\begin{figure}[H]
  \centering
  \begin{tikzpicture}
    \begin{axis}[
      ybar,
      symbolic x coords={sepolia, arbitrumSepolia, baseSepolia, zksyncSepolia},
      xtick=data,
      ylabel={Koszt \code{issueBatch(1000)} (USD)},
      bar width=10pt, legend pos=north west,
      width=0.85\textwidth, height=8cm,
    ]
      \addplot table[x=chain, y=quiet_usd] {\DataDir/fig4-basefee-scenarios.dat};
      \addplot table[x=chain, y=normal_usd] {\DataDir/fig4-basefee-scenarios.dat};
      \addplot table[x=chain, y=congested_usd] {\DataDir/fig4-basefee-scenarios.dat};
      \legend{p10 (spokój), p50 (typowo), p90 (kongestia)}
    \end{axis}
  \end{tikzpicture}
  \caption{Koszt \code{issueBatch(1000)} przy scenariuszach opłaty bazowej L1
  p10/p50/p90.}
  \label{fig:basefee}
\end{figure}

\subsection{Wrażliwość na kurs ETH/USD}
Koszt w~USD jest liniową funkcją kursu ETH/USD: $C(r) = C_{3500}\cdot r/3500$, gdzie
$C_{3500}$ to wartość bazowa raportowana przy snapshocie z~dnia 13~maja~2026~r.
($r=3500$). Tabela~\ref{tab:eth-sens} pokazuje koszt \code{issueBatch(1000)}
w~scenariuszu p50 dla trzech zakotwiczonych poziomów rynkowych.

\begin{table}[H]
  \centering
  \caption{Wrażliwość kosztu \code{issueBatch(1000)} na kurs ETH/USD
  (opłata bazowa p50).}
  \label{tab:eth-sens}
  \pgfplotstabletypeset[
    col sep=comma, string type,
    every head row/.style={before row=\toprule, after row=\midrule},
    every last row/.style={after row=\bottomrule},
    columns={chain,usd_2000,usd_3500,usd_5000},
    columns/chain/.style={column name=Sieć, column type=l},
    columns/usd_2000/.style={column name={2000 USD}, column type=r},
    columns/usd_3500/.style={column name={3500 USD}, column type=r},
    columns/usd_5000/.style={column name={5000 USD}, column type=r},
  ]{\DataDir/table-sensitivity-eth-usd.csv}
\end{table}

\subsection{Wrażliwość na cenę gazu sekwensera L2}
Cena gazu ustalana przez operatora sekwensera L2 jest niezależna od opłaty bazowej
L1. Model kosztu rozkłada się addytywnie:
\[
  C_{\text{total}} = C_{\text{exec}} + C_{\text{L1 data}},
\]
gdzie $C_{\text{exec}} = \code{gasUsed}\cdot p_{\text{L2}}$ skaluje się liniowo
z~ceną sekwensera $p_{\text{L2}}$, natomiast $C_{\text{L1 data}}$ pochodzi z~opłaty
blob na L1 i~pozostaje stała. Tabela~\ref{tab:l2price-sens} pokazuje koszt
\code{issueBatch(1000)} przy trzech mnożnikach względem snapshotu cen (Arbitrum
0{,}1~gwei, Base 0{,}005~gwei, zkSync 0{,}05~gwei). Sepolię pominięto --- jako sieć
L1 nie ma ceny sekwensera.

\begin{table}[H]
  \centering
  \caption{Wrażliwość kosztu \code{issueBatch(1000)} na cenę gazu sekwensera L2
  (mnożniki względem snapshotu cen).}
  \label{tab:l2price-sens}
  \pgfplotstabletypeset[
    col sep=comma, string type,
    every head row/.style={before row=\toprule, after row=\midrule},
    every last row/.style={after row=\bottomrule},
    columns={chain,usd_05x,usd_10x,usd_20x},
    columns/chain/.style={column name=Sieć, column type=l},
    columns/usd_05x/.style={column name={$0{,}5\times$}, column type=r},
    columns/usd_10x/.style={column name={$1\times$}, column type=r},
    columns/usd_20x/.style={column name={$2\times$}, column type=r},
  ]{\DataDir/table-sensitivity-l2-price.csv}
\end{table}

\section{Dyskusja}
\label{sec:l2-dyskusja}

{\sloppy
Wszystkie pomiary opisane w~niniejszym rozdziale wykonano przez jednego dostawcę
RPC w~jednej lokalizacji, z~$N{=}30$ powtórzeniami dla każdego pomiaru opóźnienia
oraz losowym wyborem liści przy unieważnianiu. Raportowane wyniki to średnia
$\pm\sigma$ z~trzech niezależnych przebiegów. Wszystkie parametry konfiguracyjne
(ziarno losowości, dostawca, region, kurs ETH/USD) zapisano w~pliku
\code{meta.json} towarzyszącym eksportowi danych.\par}

\subsection{Ograniczenia pomiaru}
Poniżej zebrano ograniczenia metodologiczne, których czytelnik powinien być
świadomy, interpretując wyniki. Cztery z~pierwotnie zidentyfikowanych ograniczeń
wyeliminowano w~toku rewizji eksperymentu (jednolity dostawca RPC, $N{=}30$
powtórzeń, losowy wybór liści do unieważnienia, trzy niezależne przebiegi
z~raportowaniem $\sigma$); pozostałe sześć omówiono poniżej.

\paragraph{Założenie kompresji Brotli dla Arbitrum.} Model kosztu zapisu danych na
L1 dla Arbitrum przyjmuje współczynnik kompresji Brotli równy 1{,}0, ponieważ dane
Merkle (hasze pseudolosowe) nie kompresują się znacząco. Realny calldata
o~strukturze rzadkiej kompresuje się 2--4-krotnie, co oznacza, że nasze
oszacowanie kosztu L1 dla Arbitrum jest \textbf{górnym ograniczeniem}; rzeczywiste
koszty produkcyjne mogą być proporcjonalnie niższe.

\paragraph{Niewspółmierność jednostki gazu w~zkSync Era.} Wartość \code{gasUsed}
raportowana przez zkSync (131\,463 dla \code{issueBatch} z~$n{=}1$) łączy w~jednej
jednostce koszt wykonania, dane publiczne oraz narzut bootloadera kontekstu Account
Abstraction --- nie jest to ta sama jednostka co \code{gasUsed} w~EVM. Porównywanie
surowych liczb gazu między rodzinami rollupów jest \textbf{metodologicznie
niepoprawne}; sensowne porównanie kosztów możliwe jest wyłącznie w~USD lub ETH.

\paragraph{Opóźnienie inkluzji ograniczone interwałem odpytywania.} Pomiar czasu od
wysłania do paragonu opiera się o~cykl odpytywania klienta viem
(\textasciitilde 4~s). Otrzymane wartości są zatem \textbf{górnym ograniczeniem
czasu inkluzji w~bloku}, a~nie samym czasem inkluzji. Porównanie międzysieciowe jest
dodatkowo zniekształcone interakcją czasu bloku z~interwałem odpytywania.

\paragraph{Snapshot skalarów Base SystemConfig.} Model kosztu Base (Ecotone) używa
snapshotu parametrów \code{baseFeeScalar} oraz \code{blobBaseFeeScalar} pobranego
z~konkretnego bloku L1. Po zmianie tych skalarów przez operatora Base wynikowe
koszty zmieniłyby się proporcjonalnie; data snapshotu jest udokumentowana
w~\code{meta.json} każdego eksportu.

\paragraph{Snapshot kursu ETH/USD.} Wyniki w~USD przeliczono kursem
$\text{ETH}/\text{USD}=3500$ (snapshot z~dnia 13~maja~2026~r.). Wynik dla innego
kursu otrzymuje się mnożeniem przez czynnik $r/3500$.

\paragraph{Konserwatywne ceny sekwensera L2.} Ceny gazu sekwensera ustawiono na
zaokrąglone wartości powyżej żywego snapshotu (Arbitrum 0{,}1 wobec 0{,}02~gwei;
Base 0{,}005 wobec 0{,}006; zkSync 0{,}05 wobec 0{,}0453), aby pochłonąć typowe
wahania kongestii.

\subsection{Reprodukowalność}
Wszystkie wyniki w~tym rozdziale można odtworzyć z~czystego klona repozytorium
UniVerify poleceniem \code{npm run benchmarks:all -{}- -{}-runs=3}, po wypełnieniu
pliku \code{.env} kluczem prywatnym portfela testowego oraz adresami RPC czterech
sieci testowych. Pipeline wykonuje sekwencję \code{deploy} $\to$ \code{measure}
$\to$ \code{price} $\to$ \code{aggregate} $\to$ \code{export}, regenerując wszystkie
tabele i~wykresy. Każdy artefakt powiązany jest z~plikiem \code{meta.json}
zawierającym identyfikator przebiegu, datę pomiaru, kurs ETH/USD wraz z~datą
snapshotu, ziarno losowości, dostawcę i~region RPC oraz listę sieci.

\subsection{Wnioski}
\label{sec:l2-wnioski}
Zebrane pomiary pozwalają odpowiedzieć na pytanie postawione na wstępie rozdziału:
przeniesienie rejestru do warstwy drugiej istotnie obniża koszt, lecz skala
korzyści zależy od rodziny rollupa. Przy realistycznym rozmiarze paczki $n{=}1000$
(rzędu rocznego wolumenu dyplomów dużego wydziału) koszt zarejestrowania jednego
dyplomu wynosi w~scenariuszu p50 około $3{,}2\cdot10^{-5}$~USD na~Sepolii (L1),
$1{,}8\cdot10^{-5}$~USD na~Arbitrum, $2{,}3\cdot10^{-5}$~USD na~zkSync oraz zaledwie
$9{,}0\cdot10^{-7}$~USD na~Base --- czyli około 35-krotnie taniej niż w~warstwie
pierwszej (por. tab.~\ref{tab:per-diploma}, rys.~\ref{fig:cost-per-diploma}).

Ponieważ koszt na dyplom maleje proporcjonalnie do $1/n$ (publikujemy jeden
32-bajtowy korzeń niezależnie od rozmiaru paczki), stosunki kosztów między sieciami
są niezmiennicze względem $n$; powyższa hierarchia utrzymuje się dla dowolnego
rozmiaru paczki. Przewaga Base nie wynika z~lepszej kompresji \emph{naszych}
danych --- we wszystkich sieciach publikujemy ten sam korzeń o~koszcie około
51~tys.\ gazu --- lecz z~tego, że pojedyncza transakcja L2 jest tania: rollup
amortyzuje koszt zapisu danych na~L1 pomiędzy wszystkich swoich użytkowników poprzez
transakcje blob (EIP-4844). Dla naszej transakcji składnik danych L1 stanowi znikomą
część kosztu, w~którym dominuje wykonanie po stronie sekwensera.

Wynik ten jest odporny na typową zmienność rynkową: analiza wrażliwości
(sekcja~\ref{sec:l2-wrazliwosc}) pokazuje, że nawet dwukrotne wahania kursu ETH/USD,
opłaty bazowej L1 czy ceny sekwensera L2 nie zmieniają uporządkowania sieci ani
rzędu wielkości kosztu. Z~perspektywy kosztowej Base jest zatem najkorzystniejszym
celem wdrożenia rejestru dyplomów spośród badanych sieci.
